\section{Kildestrategi og ruteplanlegging}
\label{sec:kildestrategi}

Når byens primærvannverk og rørnett er slått ut, er melkeflåtens styrke dens geografiske mobilitet. I stedet for å rense lokalt flomvann, oppretter prosjektet en «virtuell rørledning» over vei fra fire uavhengige reservekilder utenfor det rammede nedbørsfeltet.

\subsection{De fire uavhengige reservevannkildene}
\begin{enumerate}[leftmargin=1.5em, itemsep=0.25em]
    \item \textbf{Næringsmiddelindustriens dypbrønner og prosessanlegg:}
    Regionale meierier, bryggerier og slakterier disponerer egne dype grunnvannsbrønner og industrielle renseanlegg med flertrinns barrierer (ultrafiltrering, omvendt osmose og \textsc{uv}-bestråling). 
    \textit{Logistisk gevinst:} Tankbilen steriliseres i meieriets \textsc{cip}-anlegg, rygger 30 meter til fabrikkens interne rentvannsuttak, fylles i lukket sløyfe og sendes direkte ut.
    \item \textbf{Uberørte nabovannverk (Interkommunal veiforbindelse):}
    Nabokommunen 30--50 km unna henter vann fra et adskilt, uforurenset nedbørsfelt. Mange kommuner mangler fysiske overføringsrør grunnet topografi. Melkebilene fungerer som en fleksibel bypass: De kobler seg til nabovannverkets hydranter via adaptersettet og frakter vannet over kommunegrensen.
    \item \textbf{Skjermede grunnvannsbrønner i dype løsmasser:}
    Akviferer i dype sand- og gruslag påvirkes ikke av overflateflom og kloakkoverløp. Tykke sedimentlag danner et naturlig geologisk filter som sikrer mikrobiologisk trygt vann. Bilene kobles direkte til pumpehusenes ventiler.
    \item \textbf{Vannkraftmagasiner og alpine inntak:}
    Flomforurensningen skjer i lavlandet der avløp flyter over. I alpine inntaksdammer oppstrøms for bosetning finnes store mengder uberørt smelte- og regnvann. Bilene kan fylles direkte fra kraftverkenes driftsvannveier eller tømmeventiler.
\end{enumerate}

\subsection{Aksellast, flominfrastruktur og \textsc{vts}-integrasjon}
En fullastet tre- eller fireakslet melkebil har en totalvekt på 32 til 50 tonn. Under flom trues veiinfrastrukturen av erosjon og svekkede broer:
\begin{itemize}[leftmargin=1.5em, itemsep=0.2em]
    \item \textbf{Sanntidskobling mot Vegtrafikksentralen (\textsc{vts}):} Flåtestyringen integreres mot Statens vegvesens \textsc{vts}-database over stengte veier. Navigasjonssystemet ruter bilene utelukkende langs verifiserte \textsc{bk}10-korridorer (10 tonns aksellast / 50--60 tonns totalvekt).
    \item \textbf{Sikkerhetskorridorer:} Transporten ledes over høyereliggende europaveier og fylkesveier utenom flomutsatte stikkrenner. Ved usikker bæreevne på sekundære broer benyttes toakslede biler (\SI{15}{\cubic\meter}) for redusert marktrykk.
\end{itemize}
